% ===================================================================
%  ТАБЛИЦА № 13 — Гостиница. — Ресторан. — Кафе.
%  Traduction 2026 d'apres texto/io, controlee sur texto/fr, texto/en
%  et texto/es. Consigne : voir texto/ru/10-tablica-01.tex.
%
%  Deux notes, marquees toutes deux « (*) », et deux appels : celui du
%  bloc 01-2 (la chambre) et celui du bloc 09-1 (le menu). L'ordre les
%  departage.
% ===================================================================

%%K t13-noto-1 noto
\VUnotes{143.2mm}{%
\textit{\VUgras{(*) Подробности спальни см. в таблице № 6.}}%
}

%%K t13-apar-1 apar
\VUsaut{3.0mm}
\VUtitre{15.0pt}{60}{ТАБЛИЦА № 13}
\VUsaut{1.6mm}
\VUpk{10.2pt}{40}{Гостиница. --- Ресторан.}
\VUsaut{2.0mm}
\VUpk{10.2pt}{40}{Кафе.}
\VUsaut{3.4mm}

%%K t13-01-1 p
1. --- Приехав вчера вечером в Руайян, я остановился в
\textit{Гостинице Океана}, которую мне рекомендовал один приятель.

%%K t13-01-2 p
Мне дали прекрасную \VUgras{комнату} \textsuperscript{(2)} на втором
\VUgras{этаже} \textsuperscript{(1)}, где я отлично выспался (*).

%%K t13-02-1 p
2. --- Сегодня утром, выйдя из своей комнаты, куда \VUgras{горничная}
\textsuperscript{(3)} принесла мне завтрак, я зашёл в \VUgras{курительную}
\textsuperscript{(5)}, которая служит и \VUgras{читальней} \textsuperscript{(4)}, чтобы
просмотреть \VUgras{газеты} \textsuperscript{(6)}, покуривая \VUgras{папиросу}
\textsuperscript{(8)}. Кончив читать, я спустился в \VUgras{лифте} \textsuperscript{(9)}
и пошёл прогуляться по городу.

%%K t13-03-1 p
3. --- Так как я забыл спросить у гостиничного \VUgras{конторщика}
\textsuperscript{(12)}, в котором часу подают за \VUgras{общим столом}
\textsuperscript{(11)}, я вернулся рано и воспользовался случаем принять ванну
в гостиничной \VUgras{ванной} \textsuperscript{(10)}. Потом я пошёл к
\VUgras{кассе} \textsuperscript{(15)}, чтобы позвонить одному приятелю. Но
\VUgras{телефон} \textsuperscript{(13)} был занят; видя моё затруднение,
\VUgras{кассирша} \textsuperscript{(14)}, вносившая \VUgras{счёт} \textsuperscript{(17)}
в \VUgras{книгу приезжающих} \textsuperscript{(16)}, попросила \VUgras{швейцара}
\textsuperscript{(18)} послать \VUgras{рассыльного} \textsuperscript{(19)} исполнить моё
поручение на его \VUgras{велосипеде} \textsuperscript{(20)}.

%%K t13-04-1 p
4. --- Я поблагодарил её и вышел дать на чай \VUgras{носильщику}
\textsuperscript{(24)} (человеку на все руки), который привёз с вокзала на своей
\VUgras{ручной тележке} \textsuperscript{(25)} мою \VUgras{шляпную картонку}
\textsuperscript{(26)}, мой \VUgras{чемодан} \textsuperscript{(27)} и мой
\VUgras{дорожный мешок} \textsuperscript{(28)}. \VUgras{Почтальон} \textsuperscript{(21)},
подошедший как раз в ту минуту, подал мне \VUgras{письмо} \textsuperscript{(22)}.

%%K t13-05-1 p
5. --- Я задержался ещё немного на пороге, возле \VUgras{почтового ящика}
\textsuperscript{(23)}, глядя на движение по улице. \VUgras{Кучер} \textsuperscript{(30)}
\VUgras{омнибуса} \textsuperscript{(29)} грузил на свою повозку \VUgras{сундук}
\textsuperscript{(31)} одного \VUgras{постояльца} \textsuperscript{(32)}, с которым
прощался \VUgras{хозяин гостиницы} \textsuperscript{(33)}. Молодой
\VUgras{телеграфист} \textsuperscript{(35)} нёс, ничуть не торопясь,
\VUgras{депешу} \textsuperscript{(34)} одному из постояльцев;
\VUgras{турист} \textsuperscript{(36)} с \VUgras{фотографическим аппаратом}
\textsuperscript{(38)} через плечо справлялся со своим \VUgras{путеводителем}
\textsuperscript{(37)}.

%%K t13-tit-1 sub
\VUsaut{4.0mm}
\VUpk{10.2pt}{40}{Час завтрака.}
\VUsaut{3.0mm}

%%K t13-06-1 p
6. --- Перед решёткой безмятежный \VUgras{посыльный} \textsuperscript{(39)}
дремал между своими \VUgras{крючьями} \textsuperscript{(40)} и своим
\VUgras{ящиком чистильщика} \textsuperscript{(41)}, а молодая \VUgras{прачка}
\textsuperscript{(42)} с \VUgras{корзиной} \textsuperscript{(43)} белья смеялась над
важным видом \VUgras{метрдотеля} \textsuperscript{(44)}, стоявшего у двери.

%%K t13-07-1 p
7. --- Вид \VUgras{повара} \textsuperscript{(45)}, который вышел из своей
\VUgras{кухни} \textsuperscript{(47)} в \VUgras{подвале} \textsuperscript{(48)}, чтобы
послать \VUgras{поварёнка} \textsuperscript{(46)} с поручением, напомнил мне, что
близок час завтрака, и я пошёл в ресторан на другой стороне двора, где
\VUgras{автомобилист} \textsuperscript{(50)} ставил свой \VUgras{автомобиль}
\textsuperscript{(49)}, а \VUgras{механик} \textsuperscript{(52)} чинил, по указаниям
\VUgras{велосипедиста} \textsuperscript{(53)}, \VUgras{бензиновый трицикл}
\textsuperscript{(51)}, только что потерпевший аварию.

%%K t13-08-1 p
8. --- Я спросил у молодого \VUgras{грума} \textsuperscript{(54)}, который нёс
\VUgras{кружки пива} \textsuperscript{(55)}, под верандой ли общий стол, и, когда
он сказал, что да, сел за один из столиков, и мне подали превосходный
завтрак.

%%K t13-noto-2 noto
\VUnotes{143.2mm}{%
\textit{\VUgras{(*) С помощью списка блюд, данного в словаре, учитель
легко может разнообразить приведённое ниже меню.}}%
}

%%K t13-tit-2 sub
\VUsaut{4.0mm}
\VUpk{10.2pt}{40}{Отличный завтрак.}
\VUsaut{3.0mm}

%%K t13-09-1 p
9. --- Лакей принёс мне меню завтрака (*) и карту вин. Я выбрал и заказал
на \VUgras{закуску} \VUgras{креветок} с \VUgras{маслом}, а на первое ---
\VUgras{рубец по-кански}. \VUgras{Рыбы} я не брал, но попросил порцию
бараньей \VUgras{ноги} и, на \VUgras{овощ}, \VUgras{шпината} со сливками.
Потом, так как ни одно из \VUgras{сладких блюд} мне не приглянулось, я
велел подать разные десерты: \VUgras{сыр}, \VUgras{фрукты} и
\VUgras{печенье}. Я прибавил превосходного сент-эмильонского \VUgras{вина}
и, очень довольный своим завтраком, встал из-за стола, дав
\VUgras{лакею} \textsuperscript{(57)} щедрые \VUgras{чаевые} \textsuperscript{(59)}.

%%K t13-10-1 p
10. --- Видя, что я собираюсь уходить, \VUgras{управляющий} \textsuperscript{(60)}
предложил мне выйти на балкон полюбоваться великолепной панорамой
\VUgras{океана} \textsuperscript{(62)}. Вдали различались маленькие рыбачьи
\VUgras{лодки} \textsuperscript{(63)}, \VUgras{парусники} \textsuperscript{(64)} и
огромный \VUgras{пакетбот} \textsuperscript{(65)}, чей чёрный силуэт вырезался на
белых \VUgras{облаках} \textsuperscript{(67)} у \VUgras{горизонта} \textsuperscript{(66)}.
Лёгкий ветерок шевелил \VUgras{флаг} \textsuperscript{(70)}, водружённый над
\VUgras{вывеской} \textsuperscript{(69)} \VUgras{вестибюля} \textsuperscript{(68)}.

%%K t13-tit-3 sub
\VUsaut{3.0mm}
\VUpk{10.2pt}{40}{Развлечения.}
\VUsaut{3.0mm}

%%K t13-11-1 p
11. --- Зрелище было так занятно, что я решил на другой день подняться на
террасу кафе, взяв с собою \VUgras{бинокль} \textsuperscript{(71)} или
\VUgras{подзорную трубу} \textsuperscript{(72)}.

%%K t13-12-1 p
12. --- Уйдя с балкона, я пошёл в гостиничное кафе выпить \VUgras{чего-нибудь}
\textsuperscript{(73)} и сыграть партию на \VUgras{бильярде} \textsuperscript{(75)}. Я
прошёл прямо через залу кафе, где много посетителей играло в
\VUgras{шахматы} \textsuperscript{(78)}, в \VUgras{шашки} \textsuperscript{(79)}, в
\VUgras{домино} \textsuperscript{(80)}, в \VUgras{триктрак} \textsuperscript{(81)} или в
\VUgras{карты} \textsuperscript{(82)}, и поднялся прямо в \VUgras{бильярдную}
\textsuperscript{(74)}.

%%K t13-13-1 p
13. --- Кончив партию, я спустился в нижний этаж и попросил у лакея
\VUgras{конверт} \textsuperscript{(84)}, \VUgras{бумаги} \textsuperscript{(83)},
\VUgras{марку} \textsuperscript{(85)}, \VUgras{перо} \textsuperscript{(87)} и
\VUgras{чернил} \textsuperscript{(86)}, чтобы написать письмо.

%%K t13-13-2 p
Потом я подозвал \VUgras{извозчика} \textsuperscript{(92)}, чей \VUgras{экипаж}
\textsuperscript{(88)} остановился перед кафе. Я спросил его цену за час и за
конец, и, когда мы сговорились, он открыл мне \VUgras{дверцу}
\textsuperscript{(89)} и усадил внутрь. Он снова взобрался на свои \VUgras{козлы}
\textsuperscript{(91)}, бойко тронул лошадей своим \VUgras{кнутом}
\textsuperscript{(93)}, и мы покатили на восхитительную прогулку.
\VUsaut{6.0mm}
\VUfilet{22mm}
